

\begin{frame}{Litigation Pushes Voters Off the Mayoral Ballot}
\hypertarget{src:voterballot}{}
\small
Turnout is compulsory in Brazil for literate voters aged 18--69 \ns{(CF/88 art.~14 \S1)} and so is institutionally floored, but a voter can still decline to choose, casting a blank or null ballot that elects no one. With turnout floored, demobilization would appear as a rise in these non-valid votes and a fall in valid ones. The mayoral estimates below carry that sign pattern, and only the blank rise approaches significance.
\smallskip
\begin{center}
\scriptsize
\resizebox{!}{0.22\textheight}{% Auto-generated by code/03_estimation/02_iv_main.R
% Do not edit manually -- rerun the generating script to update

\begin{tabular}{lccc}
\toprule\toprule
 & \textbf{Blank vote} & \textbf{Null vote} & \textbf{Valid vote} \\
\midrule
\rowcolor{mylight}
\quad \textbf{Judicialization} & $0.003$$^{*}$ & $0.003$ & $-0.007$ \\
\rowcolor{mylight}
 & \textcolor{mygray}{(0.002)} & \textcolor{mygray}{(0.002)} & \textcolor{mygray}{(0.006)} \\
\quad {\footnotesize 2024 Mean} & {\footnotesize 0.016} & {\footnotesize 0.027} & {\footnotesize 0.792} \\
\quad {\footnotesize 2016 Mean} & {\footnotesize 0.016} & {\footnotesize 0.051} & {\footnotesize 0.789} \\[2pt]
\midrule
Municipalities ($N$) & \multicolumn{3}{c}{5,560} \\
\bottomrule\bottomrule
\end{tabular}
}
\end{center}
\smallskip
\takeaway{A one-SD shock raises the mayoral blank rate by $+\BlankPerSD$~pp and the null rate by $+\NullPerSD$~pp, cutting valid votes by $\ValidPerSDSigned$~pp. The blank rise clears the exposure-robust SE ($p=\ERBlankPbhj$), is conventional only at 10\% ($p=\BlankP$), and fails the weak-IV AR bootstrap ($p=\ARBlankP$). Null ($p=\NullP$) and valid ($p=\ValidP$) share the sign pattern without significance. In contested seats the blank rise clears 5\% ($p=\ContBlankP$). The mayoral valid rate fails the pre-trend placebo ($p=\PtRFValidP$). The council ballot is flat.}\hfill\hyperlink{app:turnoutplacebo}{\beamergotobutton{turnout placebo}}\,\hyperlink{app:councilplacebo}{\beamergotobutton{council placebo}}
\end{frame}
